\section{Study Design and Predictive Models}
\label{sec:method}

\subsection{Panels}
\label{sec:panels}

Three kinds of panel serve three purposes (Appendix Table~\ref{tab:panels}).
The \emph{heterogeneous panel} of finished public models tests how far
a relation carries across families whose training history, data, and
architecture all differ at once; it establishes the phenomena and the
range of responses, and its family and data provenance are confounded
by construction. The \emph{controlled Pythia panel} varies the source
state along known axes, size and pretraining stage, with architecture
and data held fixed, and is the development set of every predictor.
The \emph{distillation panels} vary the student, the data pool, and the
exposure with the teacher fixed.


\subsection{Predictors that were actually tested}
\label{sec:predictors}

Every predictor below was fit on the development set only, with its
standardization constants, coefficients, and keep rules committed to a
register file before any test measurement (paths in
Appendix~\ref{sec:app_register}). Parameter counts are per capability.

\paragraph{Configuration-indicator source regression (source axis).}
At each intervention configuration $q$, a density or a bit-width, a
separate linear model in the standardized inputs,
$\widehat{\Delta L}_{c}(q)=\beta_{c,q}\cdot(1,\,\tilde{\log N_0},\,\tilde L_{0,c},\,\tilde{\log D_0})$,
four coefficients per configuration (16 for four densities). It
predicts a new source state at a \emph{seen} configuration and is not a
law over the intervention axis. Its same-input comparator drops
$\tilde{\log D_0}$ (three coefficients per configuration). This is the
predictor frozen for the step-96k tests (register \texttt{v38}) and,
per bit-width, the quantization predictor throughout.

\paragraph{Shared power form (pruning, configuration axis).}
$\widehat{\Delta L}_c(\mathbf x,d)=A_c(\mathbf x)\,((1-d)/0.3)^{\gamma_c}$
with $A_c(\mathbf x)=\beta_c\cdot(1,\tilde{\log N_0},\tilde L_{0,c},\tilde{\log D_0})$,
five parameters, fit on the nine development states at
$d\in\{0.9,0.8,0.7,0.6\}$ (108 rows; register \texttt{v40}, refrozen
with the same form in \texttt{v49}). It is the headline candidate for
unseen densities.

\paragraph{Same-input alternatives for pruning.} Three forms use
exactly the inputs of the shared power form and a different strength
structure: A1, the power form with $\gamma_c$ fixed to one (four
parameters); A2, the configuration-indicator regression at the four
development densities followed by a fixed rule, linear interpolation
between adjacent predicted anchors and linear extension beyond the
boundary pair; and a continuous two-term polynomial in $s=1-d$ with
source-dependent coefficients. They test whether the power structure
itself adds anything.

\paragraph{Source-free baselines for pruning.} A strength-only curve
$A\,((1-d)/0.3)^{\gamma}$ fit on all development rows (two
parameters); the median development curve, the per-density median of
the development responses extended by the same power rule; and zero
change. They answer whether the source state is worth reading at all.

\paragraph{Quantization baselines.} For each bit-width, the per-bit
mean and median of the development responses and zero change. The
5-bit prediction of every candidate is a fixed rule, the average of its
committed 4- and 6-bit predictions, so that no 5-bit measurement enters
any fit.

\paragraph{Distillation predictors.} On the source axis, a linear model
in $(\tilde{\log N_0},\tilde L_{0,c},\tilde{\log D_0})$ (four
parameters), its no-$D_0$ variant, and per-capability constants (mean
and median of the development responses). On the pool axis, four forms
fit on the 18 endpoint checkpoints of pools U75 and U600: a constant,
$a+b\log T$, $a+b\log E$, and $a+b\log T+c\log E$; all four were
frozen and applied to the six U225 runs (register \texttt{v41}).

\subsection{Validation protocol and labels}
\label{sec:protocol}

Two axes are tested separately. On the \emph{source axis} the
intervention setting is seen and the source state is not, either
through leave-one-step-out folds on the nine-state panel or through a
frozen prospective on a new source. On the \emph{configuration axis}
the intervention setting is unseen: densities, bit-widths, or pools
outside the fitted set, with or without a new source. For the
new-source tests the order is fixed. Forms and standardization are
committed; the dense capability vector of the new source is measured,
since it is a permitted input; every candidate's predictions are
written to a versioned file and committed with the hash recorded; only
then is the compressed or distilled outcome measured. Two fitting
protocols are used in the last round: protocol A fits on all nine
development states, protocol B only on states at step 64k or below, so
that a 112k test stage is a stage extrapolation under B and a stage
interpolation under A.

Each reported comparison carries a four-part origin code (defined with
Tables~\ref{tab:pred_source}--\ref{tab:pred_config_qd}): whether the
candidate prediction was frozen before the measurement or is a
leave-one-out estimate, whether the same-input baseline and the simple
baseline were pre-specified with the candidate or computed afterwards,
and whether all pre-specified forms are reported. A frozen prediction
keeps its label whatever the outcome. When several baselines exist,
the strongest per capability is shown and that choice is marked as
retrospective. The only post-test selection in this paper is the
per-capability reading of the four distillation forms, and it is
labeled as such.

\paragraph{Which baseline answers which question.} Whether the
response is worth predicting at all: zero change and the per-capability
constant. Whether the source state helps beyond the strength curve:
strength-only and per-configuration constants against the
source-conditioned model. Whether $D_0$ adds information once $N_0$ and
$L_0$ are known: the full model against its no-$D_0$ variant on the
same split. Whether a particular strength structure helps: the shared
power form against A1, A2, and the continuous form on the same inputs.
``Same budget'' means the same admitted information; a baseline need
not use all of it.

\paragraph{Statistics.} Errors are mean absolute errors in nats per
token per capability, and improvement is baseline MAE minus candidate
MAE (positive means the candidate is better). Bootstrap intervals are
computed only on the leave-one-out panels, resampling source states as
clusters; the frozen prospectives cover one or two source states and
are reported as point values with the number of independent sources.
Repeated probes, checkpoints of one run, and stages of one source are
never counted as independent sources.
